\documentclass{article}
\usepackage{amsmath, amssymb, booktabs}

\begin{document}

\title{Solutions to Practice Exercise on the Ricardian Model}
\author{}
\date{}
\maketitle

\section*{Problem 1: Comparative Advantage and Opportunity Cost}
Two countries, \textbf{A} and \textbf{B}, produce two goods: \textbf{Wheat} and \textbf{Cloth}. The labor productivity (output per worker per hour) is given below:

\begin{table}[h]
    \centering
    \begin{tabular}{lcc}
        \toprule
        Country & Wheat (units per hour) & Cloth (units per hour) \\
        \midrule
        \textbf{A} & 10 & 5 \\
        \textbf{B} & 4 & 2 \\
        \bottomrule
    \end{tabular}
\end{table}

\subsection*{Solution:}
(a) Opportunity cost of wheat in country A: \( \frac{5}{10} = 0.5 \) units of cloth.  \\
Opportunity cost of wheat in country B: \( \frac{2}{4} = 0.5 \) units of cloth.

(b) Opportunity cost of cloth in country A: \( \frac{10}{5} = 2 \) units of wheat.  \\
Opportunity cost of cloth in country B: \( \frac{4}{2} = 2 \) units of wheat.

(c) Since the opportunity costs are the same, neither country has a comparative advantage.

\section*{Problem 2: Gains from Specialization and Trade}
Consider two countries, \textbf{X} and \textbf{Y}, producing \textbf{Cars} and \textbf{Computers}. The total labor force is \textbf{1,000 workers} in each country. The labor productivity is:

\begin{table}[h]
    \centering
    \begin{tabular}{lcc}
        \toprule
        Country & Cars per worker per year & Computers per worker per year \\
        \midrule
        \textbf{X} & 4 & 8 \\
        \textbf{Y} & 2 & 6 \\
        \bottomrule
    \end{tabular}
\end{table}

\subsection*{Solution:}
(a) Opportunity cost of 1 car in X: \( \frac{8}{4} = 2 \) computers.  \\
Opportunity cost of 1 car in Y: \( \frac{6}{2} = 3 \) computers.

(b) X specializes in cars, Y specializes in computers.  \\
Total production: \( 1000 \times 4 = 4000 \) cars and \( 1000 \times 6 = 6000 \) computers.

(c) Both countries benefit by trading at a rate between 2 and 3 computers per car.

\section*{Problem 3: Terms of Trade and Trade Gains}
Two countries, \textbf{Brazil} and \textbf{Canada}, produce \textbf{Coffee} and \textbf{Timber}. The labor productivity is as follows:

\begin{table}[h]
    \centering
    \begin{tabular}{lcc}
        \toprule
        Country & Coffee (tons per worker) & Timber (tons per worker) \\
        \midrule
        \textbf{Brazil} & 12 & 4 \\
        \textbf{Canada} & 8 & 6 \\
        \bottomrule
    \end{tabular}
\end{table}

\subsection*{Solution:}
(a) Opportunity cost of 1 ton of coffee:  \\
Brazil: \( \frac{4}{12} = 0.33 \) tons of timber.  \\
Canada: \( \frac{6}{8} = 0.75 \) tons of timber.

(b) Since Brazil has a lower opportunity cost for coffee, it should specialize in coffee, and Canada in timber.  \\
At the trade rate of 1 coffee = 0.5 timber, both countries gain as it is between their opportunity costs.

\section*{Problem 4: Effect of Technological Change on Trade Patterns}
Initially, two countries, \textbf{India} and \textbf{Germany}, produce \textbf{Steel} and \textbf{Textiles}. The productivity levels are:

\begin{table}[h]
    \centering
    \begin{tabular}{lcc}
        \toprule
        Country & Steel (tons per worker) & Textiles (meters per worker) \\
        \midrule
        \textbf{India} & 5 & 20 \\
        \textbf{Germany} & 10 & 30 \\
        \bottomrule
    \end{tabular}
\end{table}

\subsection*{Solution:}
(a) Opportunity cost of 1 ton of steel:  \\
India: \( \frac{20}{5} = 4 \) meters of textiles.  \\
Germany: \( \frac{30}{10} = 3 \) meters of textiles.

Germany has a comparative advantage in steel, and India in textiles.

(b) If Germany doubles its steel productivity, it can now produce 20 tons per worker.

New opportunity cost of steel in Germany: \( \frac{30}{20} = 1.5 \) meters of textiles.

Germany's comparative advantage in steel increases.

(c) Germany will export more steel, and India will export more textiles. The gains from trade will increase.

\end{document}
